\documentclass[11pt]{article}
\usepackage[margin=1.05in]{geometry}
\usepackage{amsmath,amssymb,amsthm,mathtools}
\usepackage{hyperref}
\usepackage{booktabs}
\hypersetup{colorlinks=true,linkcolor=blue,citecolor=blue,urlcolor=blue}

\newtheorem{thm}{Theorem}
\newtheorem{lem}[thm]{Lemma}
\newtheorem{prop}[thm]{Proposition}
\newtheorem{cor}[thm]{Corollary}
\theoremstyle{definition}
\newtheorem{defn}[thm]{Definition}
\newtheorem{rem}[thm]{Remark}
\newtheorem{hyp}[thm]{Hypothesis}

\newcommand{\At}{\widetilde{A}}
\newcommand{\bM}{\mathbf{M}}
\newcommand{\bI}{\mathbf{I}}
\newcommand{\bA}{\mathbf{A}}
\newcommand{\bB}{\mathbf{B}}
\newcommand{\bR}{\mathbf{R}}
\newcommand{\bK}{\mathbf{K}}
\newcommand{\sm}{\smallsetminus}
\newcommand{\ord}{\operatorname{ord}}
\newcommand{\Lat}{\mathcal{L}}
\newcommand{\ES}{E^{S}}

\title{The converse on the boundary divisors at full $q$,\\
and why the $A$-type divisors have none:\\
the residue of a structure constant is a difference\\
across the collision pairing}
\author{Clio}
\date{19 October 2026}

\makeatletter
%% protocol-header.tex -- PROTOCOL §2.3 page-1 block, as a TEMPLATE not an instance.
%%
%% PROTOCOL §2.3: "Every PDF must state, on its first page: author, date,
%% recipient, a title, and the commit hash of the work-in-progress push it
%% corresponds to.  A PDF with no commit hash cannot be checked against the
%% record and should be treated as a draft opinion, not a result."
%%
%% Written 2026-10-16 because the fault had recurred twice and I had twice
%% fixed only the instance:
%%   - 2026-10-14: the Q473 PDF went to Rick with no recipient and no hash.
%%   - 2026-10-14: the Q460 PDF went to Rick saying "Recipient: Robin Langer".
%%     He noticed, and took the commit hash as authority instead.
%% Both times I repaired the one file and wrote "fix the template, not just the
%% instance" in my notes.  This file is that fix.  \input it; do not retype it.
%%
%% USAGE, immediately after \maketitle:
%%   %% protocol-header.tex -- PROTOCOL §2.3 page-1 block, as a TEMPLATE not an instance.
%%
%% PROTOCOL §2.3: "Every PDF must state, on its first page: author, date,
%% recipient, a title, and the commit hash of the work-in-progress push it
%% corresponds to.  A PDF with no commit hash cannot be checked against the
%% record and should be treated as a draft opinion, not a result."
%%
%% Written 2026-10-16 because the fault had recurred twice and I had twice
%% fixed only the instance:
%%   - 2026-10-14: the Q473 PDF went to Rick with no recipient and no hash.
%%   - 2026-10-14: the Q460 PDF went to Rick saying "Recipient: Robin Langer".
%%     He noticed, and took the commit hash as authority instead.
%% Both times I repaired the one file and wrote "fix the template, not just the
%% instance" in my notes.  This file is that fix.  \input it; do not retype it.
%%
%% USAGE, immediately after \maketitle:
%%   %% protocol-header.tex -- PROTOCOL §2.3 page-1 block, as a TEMPLATE not an instance.
%%
%% PROTOCOL §2.3: "Every PDF must state, on its first page: author, date,
%% recipient, a title, and the commit hash of the work-in-progress push it
%% corresponds to.  A PDF with no commit hash cannot be checked against the
%% record and should be treated as a draft opinion, not a result."
%%
%% Written 2026-10-16 because the fault had recurred twice and I had twice
%% fixed only the instance:
%%   - 2026-10-14: the Q473 PDF went to Rick with no recipient and no hash.
%%   - 2026-10-14: the Q460 PDF went to Rick saying "Recipient: Robin Langer".
%%     He noticed, and took the commit hash as authority instead.
%% Both times I repaired the one file and wrote "fix the template, not just the
%% instance" in my notes.  This file is that fix.  \input it; do not retype it.
%%
%% USAGE, immediately after \maketitle:
%%   \input{protocol-header}
%%   \protocolheader{Rick (grandpa-rick)}{clio-vega/proofs}{60e4ce3}
%%
%% The recipient argument is MANDATORY and has no default on purpose: a default
%% is what produced "Recipient: Robin Langer" on a PDF addressed to Rick.

%% \maketitle CLEARS \@date (and \@author, \@title) in the article class, so a
%% macro called AFTER \maketitle that reads \@date prints an EMPTY row.  Caught
%% 2026-10-16 only because I ran pdftotext on page 1 instead of trusting rc=0:
%% the build was clean and the Date field was blank.  Hence: snapshot the date
%% HERE, at \input time in the preamble, where \@date is still live.
\let\protocol@date\@date

\newcommand{\protocolheader}[3]{%
  \begin{center}
  \fbox{\begin{minipage}{0.92\textwidth}
  \small
  \textbf{PROTOCOL §2.3 record.}\\[1mm]
  \begin{tabular}{@{}ll@{}}
  Author     & Clio \\
  Date       & \protocol@date \\
  Recipient  & #1 \\
  Repository & \texttt{#2} \\
  Commit     & \texttt{#3} \\
  \end{tabular}\\[1mm]
  This PDF corresponds to the pushed commit above.  Where this text and the
  commit disagree, \textbf{the commit is authoritative}.
  \end{minipage}}
  \end{center}
  \vspace{2mm}%
}

%%   \protocolheader{Rick (grandpa-rick)}{clio-vega/proofs}{60e4ce3}
%%
%% The recipient argument is MANDATORY and has no default on purpose: a default
%% is what produced "Recipient: Robin Langer" on a PDF addressed to Rick.

%% \maketitle CLEARS \@date (and \@author, \@title) in the article class, so a
%% macro called AFTER \maketitle that reads \@date prints an EMPTY row.  Caught
%% 2026-10-16 only because I ran pdftotext on page 1 instead of trusting rc=0:
%% the build was clean and the Date field was blank.  Hence: snapshot the date
%% HERE, at \input time in the preamble, where \@date is still live.
\let\protocol@date\@date

\newcommand{\protocolheader}[3]{%
  \begin{center}
  \fbox{\begin{minipage}{0.92\textwidth}
  \small
  \textbf{PROTOCOL §2.3 record.}\\[1mm]
  \begin{tabular}{@{}ll@{}}
  Author     & Clio \\
  Date       & \protocol@date \\
  Recipient  & #1 \\
  Repository & \texttt{#2} \\
  Commit     & \texttt{#3} \\
  \end{tabular}\\[1mm]
  This PDF corresponds to the pushed commit above.  Where this text and the
  commit disagree, \textbf{the commit is authoritative}.
  \end{minipage}}
  \end{center}
  \vspace{2mm}%
}

%%   \protocolheader{Rick (grandpa-rick)}{clio-vega/proofs}{60e4ce3}
%%
%% The recipient argument is MANDATORY and has no default on purpose: a default
%% is what produced "Recipient: Robin Langer" on a PDF addressed to Rick.

%% \maketitle CLEARS \@date (and \@author, \@title) in the article class, so a
%% macro called AFTER \maketitle that reads \@date prints an EMPTY row.  Caught
%% 2026-10-16 only because I ran pdftotext on page 1 instead of trusting rc=0:
%% the build was clean and the Date field was blank.  Hence: snapshot the date
%% HERE, at \input time in the preamble, where \@date is still live.
\let\protocol@date\@date

\newcommand{\protocolheader}[3]{%
  \begin{center}
  \fbox{\begin{minipage}{0.92\textwidth}
  \small
  \textbf{PROTOCOL §2.3 record.}\\[1mm]
  \begin{tabular}{@{}ll@{}}
  Author     & Clio \\
  Date       & \protocol@date \\
  Recipient  & #1 \\
  Repository & \texttt{#2} \\
  Commit     & \texttt{#3} \\
  \end{tabular}\\[1mm]
  This PDF corresponds to the pushed commit above.  Where this text and the
  commit disagree, \textbf{the commit is authoritative}.
  \end{minipage}}
  \end{center}
  \vspace{2mm}%
}

\makeatother

\begin{document}
\maketitle
\protocolheader{Rick (grandpa-rick)}{clio-vega/proofs}{1999356}

\begin{abstract}
Rick assigned me, in writing in UIDs 827 and 828, \emph{``the boundary/$A$-type
(S/A) divisors at full $q$''} --- the one family his Theorem $D^\infty$
explicitly does not settle.  The first finding is that \textbf{the assignment's
own grouping is wrong}: the two halves have different answers.

For the $A$-type divisors $\{w_i=w_j\}$ there is \textbf{no converse to be had}.
Theorem~B of \texttt{2026-10-18-atype} holds for \emph{every} unit $\Psi$ of the
$A$-type GKM ring, so no renormalisation whatever produces an $A$-type pole: the
right-hand side of any ``pole $\iff$ condition'' statement is identically false.
That is a result, not a failure, and it is recorded as such.

For the boundary divisors $S_m^{\varepsilon}=\{w_m=\varepsilon\}$, $\varepsilon^{h}=1$,
there \emph{is} a converse, and this paper proves it.  The mechanism is a single
identity.  Writing $\gamma^{\nu}_{\bM}$ for the residue at $S$ of the
$(\nu,\bM)$ entry of $\Phi^{-1}$, the boundary collapse law forces
\[
  \boxed{\ \gamma^{\nu}_{\bR\cup\{m\}}=-\,\gamma^{\nu}_{\bR}\ }
  \qquad\text{for every }\bR\subseteq[N]\sm\{m\},
\]
so the residue functional is a \emph{difference across the collision pairing}.
Feeding $f=\Psi_\chi\phi_\lambda\phi_\mu$ through it collapses $2^{N}$ terms to a
sum over
\[
  \ES(\lambda,\mu,\nu;m)
  :=\{\bR\subseteq[N]\sm\{m\}\ :\ \bR\cup\{m\}\le\nu,\ \bR\ge\lambda,\ \bR\ge\mu\},
\]
with prefactor $1-\chi(\varepsilon)$ and every term individually nonzero
(Theorem~\ref{thm:S}).  Hence: no pole if $\chi(\varepsilon)=1$ or
$\ES=\varnothing$; a pole, necessarily simple, only if $\ES\ne\varnothing$; and
--- with the quantifier where it belongs ---
\[
  \ES(\lambda,\mu,\nu;m)\ne\varnothing
  \iff
  \textit{some admissible }\chi\textit{ with }\chi(\varepsilon)\ne1
  \textit{ makes }\hat c^{\nu}_{\lambda\mu}\textit{ singular at }S_m^{\varepsilon},
\]
which is Theorem~\ref{thm:Sprime}, proved by a Vandermonde independence of the
squarefree monomials $\prod_{r\in\bR}\chi(w_r)$.  For a \emph{fixed} $\chi$ the
one remaining step is that $|\ES|\ge2$ nonzero terms do not cancel; that is the
exact analogue of what Rick needed Theorem~W and $D^\infty$ for, it is stated as
the single gap, and $|\ES|\ge2$ occurs in $38\%$ ($N=3$) to $87\%$ ($N=6$) of the
firing triples, so the census is not drawn from inside its own hypothesis class.

Two by-products.  \textbf{(i)} Motegi's order is the \emph{tail-counting} order
$\bB\le\bA\iff\#\{c\in\bB:c\ge k\}\le\#\{c\in\bA:c\ge k\}\ \forall k$
(Lemma~\ref{lem:order}); transitivity, $\bA\sm\{m\}\le\bA$, and
\textbf{complementation as an order anti-automorphism} all drop out, and the last
is what makes Motegi's inverse formula's support match the down-set.
\textbf{(ii)} $\Delta^{X}_{\bA}$ of \texttt{eq:allcard-Delta-factor}
(\texttt{2610.15344} L4257--4262) \emph{is} my diagonal entry $\At_{\bA,\bA}$, so
his all-cardinality biorthogonality gives $\Phi^{-1}$ in my own notation; I
verified it first-hand against my symmetrisation engine, $336/336$.
\end{abstract}

\tableofcontents

%=====================================================================
\section{Scope determination, done first}
\label{sec:scope}

The brief forbade assuming the target was open, and required a paragraph on
\emph{what exactly is open} before any attempt.  Here it is.

\begin{center}
\begin{tabular}{llll}
\toprule
family & root & untwisted $c^{\nu}_{\lambda\mu}$ & twisted $\hat c^{\nu}_{\lambda\mu}$\\
\midrule
pair $\{w_aw_b=1\}$ & $e_a+e_b$ & poles $\iff E\ne\varnothing$ (Rick, $D^\infty$)
   & none iff \textup{(COB)} \\
boundary $\{w_m^{h}=1\}$ & $h\,e_m$ & \textbf{none} & \textbf{one-directional only}\\
$A$-type $\{w_i=w_j\}$ & $e_i-e_j$ & none & \textbf{none, for every unit $\Psi$}\\
$(Q)$ $\{1-q\cdot\text{mono}\}$ & --- & none & none\\
\bottomrule
\end{tabular}
\end{center}

\noindent
Row 3 is Theorem~B of \texttt{2026-10-18-atype-divisors-and-the-gkm-lattice.tex};
row 4 is Theorem~\texttt{thm:Q} of \texttt{2026-10-18-c2}; the untwisted entry of
row 2 is \texttt{thm:SRC} $+$ \texttt{prop:short} of \texttt{2026-10-17-c2}; the
twisted entry of row 2 is \texttt{thm:C}(2) there.

Two consequences, and they point in opposite directions.

\textbf{The $A$ half is closed and admits no converse.}  Theorem~B of
\texttt{2026-10-18-atype} is stated for an arbitrary unit $\Psi$ of the $A$-type
GKM ring $\mathcal{G}_{ij}$, not merely for a point-indexed product; and
\texttt{lem:unit} there puts $\Psi_\chi$ in $\mathcal{G}_{ij}^{\times}$ for every
$ij$-admissible $\chi$.  So \emph{no} renormalisation produces an $A$-type pole.
A statement of the form ``pole at $\{w_i=w_j\}$ $\iff$ $E^{A}\ne\varnothing$''
therefore holds with $E^{A}$ identically empty, which is to say it has no
content.  \S\ref{sec:Ahalf} records this with its mechanism.

\textbf{The $S$ half is open in exactly one place.}  \texttt{thm:C}(2) says:
regular if $\chi(\varepsilon)=1$; and if $\chi(\varepsilon)\ne1$, regular
\emph{only when} $(\phi_\lambda\phi_\mu)(\bM^{-})|_{S}=0$ for all $\bM\ni m$.
Measured against Rick's pair-divisor theorem it is missing two things: it is a
per-$(\lambda,\mu)$ statement (``\emph{every} $\hat c^{\nu}$ is regular''), not a
per-$\nu$ one; and the non-vanishing direction is not made combinatorial.  So the
target of this paper is \textbf{the per-$\nu$ converse on the boundary family at
full $q$}.

\begin{rem}[the arm-3 detector fired, and where]
The brief's not-well-posed arm came with a detector: \emph{before measuring,
check each candidate divisor is a root hyperplane of the relevant type}.  It
fires, and it fires on the $S$ half rather than the $A$ half.  In type $B$ the
boundary root is $e_m$ and the locus is $\{w_m=1\}$: the second component
$\{w_m=-1\}$ of type $C$'s $\{w_m^{2}=1\}$ is \textbf{not} a root hyperplane of
$B_N$, and the measurement in \texttt{2026-10-19-typeB-gauge-theorem.tex}
(\S``Beyond the root system'') is that the collapse law itself fails there
($0/9$ at $r\ge1$, against $6/6$ at $r=0$ where there is no head slot).  So
every statement below is quantified over $\varepsilon^{h}=1$ --- two components
in type $C$, one in type $B$ --- and the type-$B$ half of the $\varepsilon=-1$
story does not exist rather than being unproved.
\end{rem}

%=====================================================================
\section{Setup}
\label{sec:setup}

Fix $N\ge1$ and $X\in\{B,C\}$, $h=h_X$ ($2$ for $C$, $1$ for $B$).  Let
$F:=\mathbb{C}(q,w_1,\dots,w_N)$, $\mathcal{P}:=2^{[N]}$ the $T$-fixed points of
$LG(N,2N)$, and $\mathcal{A}:=F^{\mathcal{P}}$ the GKM algebra under pointwise
operations.  For $\bI\in\mathcal{P}$ let
$\phi_{\bI}(\bM):=\At^{X}_{\bM,\bI}$, Motegi's normalised type-$X$ weight
function at $u=w_{\bM}$ (\texttt{2610.15344}; the engine transcribes
\texttt{Motherfunction} L473--487 and the normalisers L1057--1062, L1361--1367).
Write $\Phi:=(\phi_{\bI}(\bM))_{\bM,\bI}$.

We use three facts from Motegi, all deep-read:
\begin{description}
\item[(M3) Triangularity (Thm 2.6(i), L603--620; order defined L582--589).]
  $\At_{\bM,\bI}=0$ unless $\bM\ge\bI$, where $\bM\ge\bI$ means
  $|\bM|\ge|\bI|$ and $m_{s-r+j}\ge i_j$ for $j=1,\dots,r$
  ($s=|\bM|$, $r=|\bI|$, both sets listed increasingly).
\item[(M4) The diagonal (\texttt{normalizationWC} L1078, \texttt{normalizationWB}
  L1386).]  $\At_{\bI,\bI}$ is a product over the positive roots, with factors
  $\prod_{j\le k}(1-w_{i_j}^{-1}w_{i_k}^{-1})$ over $\bI$ --- note $j\le k$, so
  the boundary term $(1-w_{i}^{-h})$ occurs --- together with $A$-type factors
  and $q$-factors $(1-q\cdot\text{monomial})$.
\item[{(M5) $\At_{\bM,\bI}\in R:=\mathbb{C}(q)[w^{\pm}]$ (Prop.\ L3302--3309 type
  $C$, L3722--3729 type $B$).}]
\end{description}

By (M3)+(M4), $\{\phi_{\bI}\}$ is an $F$-basis of $\mathcal{A}$; the structure
constants are $\phi_\lambda\phi_\mu=\sum_\nu c^{\nu}_{\lambda\mu}\phi_\nu$, and
$\ell_\nu\in\mathcal{A}^{*}$ denotes the $\nu$-th coordinate functional, so that
$\ell_\nu(f)=\sum_{\bM}(\Phi^{-1})_{\nu\bM}f(\bM)$.

Fix $m\in[N]$ and $\varepsilon$ with $\varepsilon^{h}=1$.  Let
$S=S_m^{\varepsilon}:=\{w_m=\varepsilon\}$, with local ring
$\mathcal{O}\subset F$, uniformiser $t:=w_m-\varepsilon$ and residue field $k$;
$\bar x$ denotes the image of $x\in\mathcal{O}$ in $k$.  For $\bM\ni m$ write
$\bM^{-}:=\bM\sm\{m\}$.  Note $k\cong\mathbb{C}(q)(w_j:j\ne m)$.

\begin{description}
\item[(BC) Boundary collapse.]  For every $\bI$ and every $\bM\ni m$,
  $\ \At^{X}_{\bM,\bI}\big|_{S}=\At^{X}_{\bM^{-},\bI}\big|_{S}$: the boundary
  collision multiplier is exactly $1$.
\item[(BD) Boundary diagonal.]  $\ord_S\At^{X}_{\bI,\bI}=[\,m\in\bI\,]$ for
  every $\bI$, at a point of $S$ generic in the sense of
  Remark~\ref{rem:generic}.
\end{description}

\noindent
(BC) in type $C$ is \texttt{thm:SRC} of \texttt{2026-10-17-c2} (proved there in
four steps, with a plant at $\varepsilon=2$ firing $24/24$); in type $B$ it is
the gauge theorem of \texttt{2026-10-19-typeB} at $\alpha=e_m$, where
$\rho^{\alpha}=\psi_B(w_m)^{-1}$ and $\psi_B(v)=(1-qv)/(1-q)$ equals $1$ on
$\{w_m=1\}$.  (BD) in type $C$ is the count inside \texttt{prop:short} of
\texttt{2026-10-17-c2}; in type $B$ it is \textbf{measured here and not proved}
($34/34$, \S\ref{sec:verif}), and every statement using it in type $B$ is graded
accordingly in \S\ref{sec:gaps}.

\begin{defn}[admissible renormalisation]\label{def:adm}
$\chi\in\mathbb{C}(q)(v)^{\times}$ is \textbf{$S$-admissible} if $\chi(w_j)$ is
regular and non-vanishing at the generic point of $S$ for every $j\in[N]$; then
$\Psi_\chi(\bM):=\prod_{m'\in\bM}\chi(w_{m'})$ lies in
$(\mathcal{O}^{\times})^{\mathcal{P}}$.  Put $\kappa:=\chi(\varepsilon)$,
$\hat\phi_{\bI}:=\Psi_\chi\phi_{\bI}$, and let $\hat c^{\nu}_{\lambda\mu}$ be the
structure constants of $\{\hat\phi_{\bI}\}$.
\end{defn}

\begin{lem}[the twisted constants are coordinates of one function]\label{lem:red}
$\hat c^{\nu}_{\lambda\mu}=\ell_\nu(f)$ with $f:=\Psi_\chi\,\phi_\lambda\phi_\mu$.
\end{lem}
\begin{proof}
$\Psi:=\Psi_\chi$ is invertible in $\mathcal{A}$, so
$\hat\phi_\lambda\hat\phi_\mu=\sum\hat c^{\nu}\hat\phi_\nu$ is equivalent to
$\Psi^{2}\phi_\lambda\phi_\mu=\Psi\sum\hat c^{\nu}\phi_\nu$, i.e.\ to
$\Psi\phi_\lambda\phi_\mu=\sum\hat c^{\nu}\phi_\nu$.  (This is the first step of
\texttt{thm:B} of \texttt{2026-10-17-c2}.)
\end{proof}

\begin{rem}[what ``generic'' has to mean, and the control that taught me]
\label{rem:generic}
(BD) is \emph{not} a statement about $q$ alone.  The positive-root product
$\At_{\bI,\bI}$ contains factors $(1-qw_mw_k)$ and $(1-qw_k/w_m)$; at
$w_m=\varepsilon$ these are $(1-q\varepsilon w_k)$ and
$(1-q\varepsilon^{-1}w_k)$, which vanish on the \emph{hypersurfaces}
$w_k=\pm q^{-1}$ in the remaining variables.  My first $N=3$ census drew the
random point $w_3=7/3$ with $q=3/7$, landed on exactly that hypersurface, and
reported $\ord_S\det\Phi=8$ against the predicted $4$ --- a confident, legible,
wholly false number, caught only because I had required the count as a control.
So ``generic'' here means: outside the finitely many hypersurfaces on which a
root factor other than $(1-w_m^{-h})$ vanishes identically on $S$.  All
computations below sample until (BD) holds for every $\bI$, verified against an
\emph{independent} closed form for the diagonal (the root product
\texttt{diag\_rootproduct}, not the symmetrisation sum), and assert it.
\end{rem}

%=====================================================================
\section{Motegi's order is the tail-counting order}
\label{sec:order}

\begin{lem}\label{lem:order}
For $\bA,\bB\subseteq[N]$ write $\bB_{\ge k}:=\#\{c\in\bB:c\ge k\}$.  Then
\[
  \bB\le\bA
  \iff
  \bB_{\ge k}\le\bA_{\ge k}\quad\text{for every }k\in[N].
\]
\end{lem}

\begin{proof}
Write $\bA=\{a_1<\dots<a_s\}$, $\bB=\{b_1<\dots<b_r\}$; the relation $\bB\le\bA$
is ``$r\le s$ and $a_{s-r+j}\ge b_j$ for $j=1,\dots,r$''.

($\Leftarrow$)  Taking $k=1$ gives $r\le s$.  Fix $j\le r$ and put $k=b_j$.  The
elements $b_j,\dots,b_r$ are all $\ge k$, so $\bB_{\ge k}\ge r-j+1$, hence
$\bA_{\ge k}\ge r-j+1$, i.e.\ at least $r-j+1$ elements of $\bA$ are $\ge b_j$;
the $(r-j+1)$-th largest of these is $a_{s-(r-j+1)+1}=a_{s-r+j}$, so
$a_{s-r+j}\ge b_j$.

($\Rightarrow$)  Fix $k$ and put $i:=\bB_{\ge k}$.  If $i=0$ there is nothing to
prove.  Otherwise $b_{r-i+1}\ge k$, and taking $j=r-i+1$ in the hypothesis gives
$a_{s-i+1}=a_{s-r+j}\ge b_{r-i+1}\ge k$; hence $a_{s-i+1},\dots,a_s$ are all
$\ge k$ and $\bA_{\ge k}\ge i$.
\end{proof}

\begin{cor}\label{cor:order}
\begin{enumerate}
\item $\le$ is a partial order on $\mathcal{P}$.
\item $\bA\sm\{m\}\le\bA$ for every $\bA\ni m$; hence every down-set
      $\downarrow\!\nu:=\{\bM:\bM\le\nu\}$ is closed under $\bM\mapsto\bM^{-}$.
\item \textup{(complementation is an order anti-automorphism)}
      $\bB\le\bA\iff\bA^{\mathrm c}\le\bB^{\mathrm c}$.
\end{enumerate}
\end{cor}

\begin{proof}
(1) The displayed criterion is a conjunction of inequalities between integers, so
it is reflexive and transitive; and $\bB\le\bA\le\bB$ forces $\bB_{\ge k}=\bA_{\ge k}$
for all $k$, whence $\bA=\bB$ (the counting function determines the set:
$k\in\bA\iff\bA_{\ge k}>\bA_{\ge k+1}$).
(2) Deleting an element decreases every $\bA_{\ge k}$ by $0$ or $1$.
(3) $(\bA^{\mathrm c})_{\ge k}=(N-k+1)-\bA_{\ge k}$, so
$\bA^{\mathrm c}_{\ge k}\le\bB^{\mathrm c}_{\ge k}$ for all $k$ is literally
$\bB_{\ge k}\le\bA_{\ge k}$ for all $k$.
\end{proof}

\begin{rem}
Corollary~\ref{cor:order}(2) is the closure my review of 2026-10-19 flagged as
``the first thing to check'' when I proposed to transport Rick's down-set
argument to the $S$ and $A$ divisors.  It holds.  Corollary~\ref{cor:order}(3) is
the fact that makes the support of Motegi's inverse formula
(Lemma~\ref{lem:inverse}) agree with the triangularity of $\Phi^{-1}$ --- the two
mechanisms give the same support precisely because complementation reverses the
order.  Verified exhaustively: $21844/21844$ pairs $(\bM,\nu)$, all $N\le7$.
\end{rem}

%=====================================================================
\section{The residue functional, and the identity that drives everything}
\label{sec:res}

Throughout this section $\chi$ is $S$-admissible and $\nu\in\mathcal{P}$ is fixed.

\begin{prop}[boundary lattice lemma; \texttt{prop:short}]\label{prop:lat}
Assume \textup{(BC)} and \textup{(BD)}.  Then
\[
  \Lat:=\operatorname{span}_{\mathcal{O}}\{\phi_{\bI}:\bI\in\mathcal{P}\}
  \;=\;
  W:=\{f\in\mathcal{O}^{\mathcal{P}}\ :\ f(\bM)\equiv f(\bM^{-})\!\!\pmod t
       \ \ \forall\,\bM\ni m\},
\]
and $t\,\mathcal{O}^{\mathcal{P}}\subseteq\Lat$.
\end{prop}

\begin{proof}
$\Lat\subseteq W$: the $\phi_{\bI}$ lie in $\mathcal{O}^{\mathcal{P}}$ by (M5)
and satisfy the congruence by (BC); $W$ is an $\mathcal{O}$-module.
Colength of $W$: the defect map
$D:\mathcal{O}^{\mathcal{P}}\to k^{\{\bM\ni m\}}$,
$D_{\bM}(f):=\overline{f(\bM)}-\overline{f(\bM^{-})}$, is $\mathcal{O}$-linear with
kernel $W$, and is surjective --- for $\bM_0\ni m$ the indicator of $\bM_0$ has
$D_{\bM}=\delta_{\bM,\bM_0}$, since $\bM^{-}=\bM_0$ is impossible for
$m\in\bM_0$.  So $\operatorname{length}(\mathcal{O}^{\mathcal{P}}/W)=2^{N-1}$.
Colength of $\Lat$: over the DVR $\mathcal{O}$ it is
$\ord_S\det\Phi=\sum_{\bI}\ord_S\At_{\bI,\bI}$ (by (M3) and
Corollary~\ref{cor:order}(1): a permutation $\sigma$ with $\sigma(\bM)\le\bM$ for
all $\bM$ is the identity, so $\det\Phi=\prod_{\bI}\At_{\bI,\bI}$), which is
$\#\{\bI\ni m\}=2^{N-1}$ by (BD).  Nested submodules of equal finite colength
coincide.  Finally $t\,\mathcal{O}^{\mathcal{P}}\subseteq W$ because
$\overline{tg(\bM)}=0=\overline{tg(\bM^{-})}$.
\end{proof}

\begin{lem}[the residue functional]\label{lem:resfun}
$t\,(\Phi^{-1})_{\nu\bM}\in\mathcal{O}$ for all $\bM$.  Hence, setting
$\gamma^{\nu}_{\bM}:=\overline{t\,(\Phi^{-1})_{\nu\bM}}\in k$,
\[
  \operatorname{Res}_\nu(f):=\overline{t\,\ell_\nu(f)}
  =\sum_{\bM\in\mathcal{P}}\gamma^{\nu}_{\bM}\,\overline{f(\bM)}
  \qquad\text{for every }f\in\mathcal{O}^{\mathcal{P}},
\]
and $\ell_\nu(f)\in\mathcal{O}\iff\operatorname{Res}_\nu(f)=0$.
\end{lem}

\begin{proof}
By Proposition~\ref{prop:lat}, $t\,\mathcal{O}^{\mathcal{P}}\subseteq\Lat$, so
for every $g\in\mathcal{O}^{\mathcal{P}}$ all coordinates $\ell_\bI(tg)$ lie in
$\mathcal{O}$; taking $g$ the indicator of $\bM$ gives
$t(\Phi^{-1})_{\bI\bM}\in\mathcal{O}$.  Consequently $\ord_S\ell_\nu(f)\ge-1$ for
$f\in\mathcal{O}^{\mathcal{P}}$, so $t\ell_\nu(f)\in\mathcal{O}$ and its class is
$\sum_\bM\overline{t(\Phi^{-1})_{\nu\bM}}\,\overline{f(\bM)}$.  The last claim is
the statement that a pole of order $\le1$ vanishes iff its residue does.
\end{proof}

\begin{lem}[$\gamma$ is a difference across the collision pairing]\label{lem:gamma}
For every $\bR\subseteq[N]\sm\{m\}$,
\[
  \gamma^{\nu}_{\bR\cup\{m\}}=-\,\gamma^{\nu}_{\bR}.
\]
Moreover $\gamma^{\nu}_{\bM}=0$ unless $\bM\le\nu$; hence $\gamma^{\nu}_{\bR}=0$
unless $\bR\cup\{m\}\le\nu$.
\end{lem}

\begin{proof}
\emph{The identity.}  $\operatorname{Res}_\nu$ vanishes on $\Lat$: if $f\in\Lat$
then $\ell_\nu(f)\in\mathcal{O}$, so $t\ell_\nu(f)\in t\mathcal{O}$.  By
Proposition~\ref{prop:lat}, $\Lat=W$, and the reduction map
$W\to k^{\mathcal{P}}$ has image exactly
$\mathcal{B}:=\{g\in k^{\mathcal{P}}:g(\bM)=g(\bM^{-})\ \forall\bM\ni m\}$: it
lands there by definition of $W$, and it is onto because given $g\in\mathcal{B}$
one may lift $g(\bR)$ arbitrarily for each $\bR\not\ni m$ and then \emph{define}
$f(\bR\cup\{m\}):=f(\bR)$, obtaining $f\in W$ with $\bar f=g$.  Therefore
\[
  \sum_{\bM}\gamma^{\nu}_{\bM}\,g(\bM)=0\qquad\text{for every }g\in\mathcal{B},
\]
and $\mathcal{B}$ is free on the coordinates indexed by $\{\bR:m\notin\bR\}$.
Taking $g$ the element of $\mathcal{B}$ with
$g(\bR_0)=g(\bR_0\cup\{m\})=1$ and all other coordinates $0$ gives
$\gamma^{\nu}_{\bR_0}+\gamma^{\nu}_{\bR_0\cup\{m\}}=0$.

\emph{The support.}  By (M3), $\Phi_{\bM\bI}=0$ unless $\bI\le\bM$, i.e.\ $\Phi$
is triangular for any linear extension of the partial order
(Corollary~\ref{cor:order}(1)); the inverse of a triangular matrix is triangular
for the same extension, so $(\Phi^{-1})_{\nu\bM}=0$ unless $\bM\le\nu$.  The last
sentence then follows from the identity, since $\gamma^\nu_\bR\ne0$ forces
$\gamma^\nu_{\bR\cup\{m\}}\ne0$ and hence $\bR\cup\{m\}\le\nu$.
\end{proof}

\begin{rem}[why this is the right object]
Lemma~\ref{lem:gamma} is where the boundary family differs structurally from the
pair family.  There the collision multiplier is $g\ne1$ and the analogous
coefficients satisfy $\gamma_{\bM}=-g\,\gamma_{\bM''}$; here it is exactly $1$,
so $\operatorname{Res}_\nu$ is a \emph{plain} difference across the pairing
$\bR\leftrightarrow\bR\cup\{m\}$.  That is what lets the $2^{N}$-term functional
collapse to a sum over a combinatorially described set without any asymptotic
input.
\end{rem}

%=====================================================================
\section{Motegi's inverse, and the one unproved input}
\label{sec:inv}

\begin{lem}[$\Phi^{-1}$, from all-cardinality biorthogonality]\label{lem:inverse}
With $|\lambda(\bA)|:=\sum_{a\in\bA}a$,
\[
  (\Phi^{-1})_{\nu\bM}
  =q^{\,|\lambda(\nu)|-|\lambda(\bM)|}\,
   \frac{\At_{\bM^{\mathrm c},\nu^{\mathrm c}}(w,q^{-1})}
        {\At_{\bM^{\mathrm c},\bM^{\mathrm c}}(w,q^{-1})\;
         \At_{\bM,\bM}(w,q)} .
\]
\end{lem}

\begin{proof}
Motegi's \texttt{thm:BC-all-cardinality-biorthogonality} (L4365--4395) reads
\[
  \sum_{\bI\subseteq[N]}q^{\,|\lambda(\bI)|-|\lambda(\mathbf L)|}
  \frac{\widetilde{WX}_{N,N,|\bI|}(w_{\bM}|w|\bI)}{\Delta^{X}_{\bM}(w,q)}
  \cdot
  \frac{\widetilde{WX}_{N,N,N-|\bI|}
        (w_{\mathbf L^{\mathrm c}}|w|\bI^{\mathrm c})\big|_{q\mapsto q^{-1}}}
       {\Delta^{X}_{\mathbf L^{\mathrm c}}(w,q^{-1})}
  =\delta_{\bM,\mathbf L},
\]
i.e.\ $\mathsf A^{X}(\mathsf A^{X})^{\sharp}=\mathrm{Id}$ where
$\mathsf A^{X}_{\bM,\bI}=\widetilde{WX}_{N,N,|\bI|}(w_\bM|w|\bI)/\Delta^{X}_{\bM}$
and $(\mathsf A^{X})^{\sharp}_{\bI,\mathbf L}
=q^{|\lambda(\bI)|-|\lambda(\mathbf L)|}\mathsf A^{X}_{\mathbf L^{\mathrm c},\bI^{\mathrm c}}(w,q^{-1})$.
Two identifications:
\begin{enumerate}
\item $\widetilde{WX}_{N,N,r}(w_{\bA}|w|\bI)=\At_{\bA,\bI}$.  Motegi pads
  $w_\bA$ with $N-|\bA|$ entries equal to $1$ and takes $n=N$ (L4250--4253); by
  iterated stability (type $C$ L2822--2828, type $B$ L3602--3608) the padded
  value at $n=N$ equals the unpadded one at $n=|\bA|$ when $|\bA|\ge r$, and both
  vanish when $|\bA|<r$ (L2843, L3620--3626).
\item $\Delta^{X}_{\bA}(w,q)=\At_{\bA,\bA}$.  This is
  \texttt{eq:allcard-Delta-factor}, L4257--4262, \emph{verbatim}:
  $\Delta^{X}_{\bA}:=\widetilde{WX}_{N,N,|\bA|}(w_\bA|w|\bA)$.
\end{enumerate}
Hence $\mathsf A^{X}=\operatorname{diag}(\At_{\bM,\bM})^{-1}\Phi$, so
$\Phi^{-1}=(\mathsf A^{X})^{\sharp}\operatorname{diag}(\At_{\bM,\bM})^{-1}$,
which is the display.
\end{proof}

\begin{rem}[gate 0b of the brief, answered, including its arm 3]\label{rem:gate}
The brief required me to verify this dependency first-hand before using it, and
pre-registered \textbf{P2}: \emph{does $\widehat\Phi$ of
\texttt{lem:allcard-Phihat-inverse} specialise to my $\At$-restriction at type
$A$?}  \textbf{Arm 3 fired: $\widehat\Phi$ is a different object.}  By
\texttt{eq:allcard-Phihat-def} (L4287--4297) $\widehat\Phi$ is the \emph{type-$A$},
fixed-cardinality matrix $\Phi_{\bK,\bI}/\Phi_{\bK,\bK}$, block diagonal in
cardinality with unit diagonal --- one of the two factors of
$\mathsf A^{X}=\widehat{\mathrm{Br}}{}^{X}\widehat\Phi$.  It is not a
restriction of my type-$B/C$ $\At$ and I did not use it.  What \emph{is} mine, up
to the diagonal row scaling in Lemma~\ref{lem:inverse}, is the \textbf{parent}
theorem's matrix $\mathsf A^{X}$, and I verified the parent theorem first-hand
against my own symmetrisation engine: $\mathbf{336/336}$ pairs $(\bM,\mathbf L)$,
$N=1,2,3$, both types, two random rational $(q,w)$ points each
(\S\ref{sec:verif}).  That is a \emph{different mechanism} from Motegi's proof
(his is branching plus RTV stable-envelope uniqueness; mine is evaluating the
symmetrisation formula), so it is corroboration rather than a second reading of
one engine --- and it discharges, from my side and at $N\le3$, Gap 1 of Rick's
$D^\infty$ document.
\end{rem}

\begin{hyp}[$L^{S}$: sharp triangularity after boundary restriction]\label{hyp:L}
If $\bA\ge\bB$ and $m\notin\bA$ then $\At_{\bA,\bB}\big|_{S}\ne0$ in $k$.
\end{hyp}

\noindent
Hypothesis~\ref{hyp:L} is the boundary analogue of Rick's Lemma~L (his \S2bis),
which gives the same sharpness without the restriction, by a degeneration
$w_c=T^{-e_c}$, $T\to0^{+}$, at $q=0$.  That degeneration cannot be run with
$w_m$ pinned at $\varepsilon$, so Hypothesis~\ref{hyp:L} is \textbf{not proved
here}.  It is measured: $699/699$ instances at $N\le4$, both types, every $m$,
every admissible $\varepsilon$, with the plant ($\bA\not\ge\bB$, where the value
\emph{must} vanish) firing $1179/1179$ (\S\ref{sec:verif}).  Every result below
that uses it says so.

\begin{lem}[$\gamma$ is nonzero on its combinatorial support]\label{lem:gnz}
Assume Hypothesis~\ref{hyp:L} and \textup{(BD)}.  If $\bM\le\nu$ and $m\in\bM$
then $\gamma^{\nu}_{\bM}\ne0$.  Equivalently, by Lemma~\ref{lem:gamma}:
$\gamma^{\nu}_{\bR}\ne0$ if and only if $\bR\cup\{m\}\le\nu$.
\end{lem}

\begin{proof}
Apply Lemma~\ref{lem:inverse}.  By (BD), $\ord_S\At_{\bM,\bM}(w,q)=1$ since
$m\in\bM$, and $\ord_S\At_{\bM^{\mathrm c},\bM^{\mathrm c}}(w,q^{-1})=0$ since
$m\notin\bM^{\mathrm c}$ --- note (BD) is stable under $q\mapsto q^{-1}$, which is
an automorphism of $\mathbb{C}(q)$ fixing the notion of genericity.  By
Corollary~\ref{cor:order}(3), $\bM\le\nu$ gives
$\nu^{\mathrm c}\le\bM^{\mathrm c}$, so by Hypothesis~\ref{hyp:L} applied with
$\bA=\bM^{\mathrm c}$ (which omits $m$) and $\bB=\nu^{\mathrm c}$, the numerator
restricts to a nonzero element of $k$.  Hence
$\ord_S(\Phi^{-1})_{\nu\bM}=-1$ exactly, and
$\gamma^{\nu}_{\bM}=\overline{t(\Phi^{-1})_{\nu\bM}}\ne0$.  The ``only if'' half
is Lemma~\ref{lem:gamma}.
\end{proof}

%=====================================================================
\section{Theorem S: the residue of a boundary structure constant}
\label{sec:S}

\begin{defn}\label{def:ES}
For $\lambda,\mu,\nu\in\mathcal{P}$ and $m\in[N]$ put
\[
  \ES(\lambda,\mu,\nu;m)
  :=\bigl\{\bR\subseteq[N]\sm\{m\}\ :\
      \bR\cup\{m\}\le\nu,\ \ \bR\ge\lambda,\ \ \bR\ge\mu\bigr\}.
\]
\end{defn}

\begin{thm}[the residue formula]\label{thm:S}
Let $\chi$ be $S$-admissible, $\kappa=\chi(\varepsilon)$, and
$f=\Psi_\chi\phi_\lambda\phi_\mu$.  Then
\[
  \operatorname{Res}_\nu(f)
  =(1-\kappa)\sum_{\bR\in \ES(\lambda,\mu,\nu;m)}
     \gamma^{\nu}_{\bR}\;
     \overline{\Psi_\chi(\bR)}\;\overline{\phi_\lambda(\bR)\phi_\mu(\bR)} .
\]
Assuming Hypothesis~\ref{hyp:L}, every summand is nonzero in $k$.
Consequently:
\begin{enumerate}
\item \textup{(no pole)} if $\kappa=1$ or $\ES(\lambda,\mu,\nu;m)=\varnothing$,
      then $\hat c^{\nu}_{\lambda\mu}\in\mathcal{O}$;
\item \textup{(necessity, and simplicity)} if $\hat c^{\nu}_{\lambda\mu}\notin\mathcal{O}$
      then $\kappa\ne1$ and $\ES(\lambda,\mu,\nu;m)\ne\varnothing$, and the pole is
      simple;
\item \textup{(sufficiency when the sum is a single term)} if $\kappa\ne1$ and
      $|\ES(\lambda,\mu,\nu;m)|=1$ then, assuming Hypothesis~\ref{hyp:L},
      $\hat c^{\nu}_{\lambda\mu}$ has a simple pole at $S_m^{\varepsilon}$.
\end{enumerate}
Items (1) and (2) do not use Hypothesis~\ref{hyp:L}.
\end{thm}

\begin{proof}
By Lemma~\ref{lem:red}, $\hat c^{\nu}_{\lambda\mu}=\ell_\nu(f)$, and
$f\in\mathcal{O}^{\mathcal{P}}$ because $\Psi_\chi$ is a unit at $S$ and
$\phi_\lambda,\phi_\mu\in\mathcal{O}^{\mathcal{P}}$ by (M5).  By
Lemma~\ref{lem:resfun},
$\operatorname{Res}_\nu(f)=\sum_{\bM}\gamma^{\nu}_{\bM}\overline{f(\bM)}$.  Group
the $2^{N}$ terms into the $2^{N-1}$ pairs $\{\bR,\bR\cup\{m\}\}$ and apply
Lemma~\ref{lem:gamma}:
\[
  \operatorname{Res}_\nu(f)
  =\sum_{\bR\not\ni m}\gamma^{\nu}_{\bR}
     \Bigl(\overline{f(\bR)}-\overline{f(\bR\cup\{m\})}\Bigr).
\]
Now compute $\overline{f(\bR\cup\{m\})}$.  By (BC),
$\overline{\phi_\lambda(\bR\cup\{m\})}=\overline{\phi_\lambda(\bR)}$ and likewise
for $\mu$; and
$\overline{\Psi_\chi(\bR\cup\{m\})}
=\overline{\chi(w_m)}\cdot\overline{\Psi_\chi(\bR)}=\kappa\,\overline{\Psi_\chi(\bR)}$.
So $\overline{f(\bR\cup\{m\})}=\kappa\,\overline{f(\bR)}$ and
\[
  \operatorname{Res}_\nu(f)
  =(1-\kappa)\sum_{\bR\not\ni m}\gamma^{\nu}_{\bR}\,\overline{f(\bR)} .
\]
Two vanishing conditions cut this down to $\ES$.  First, by
Lemma~\ref{lem:gamma}, $\gamma^{\nu}_{\bR}=0$ unless $\bR\cup\{m\}\le\nu$.
Second, by (M3), $\phi_\lambda(\bR)=0$ identically unless $\bR\ge\lambda$, and
likewise for $\mu$; so $\overline{f(\bR)}=0$ unless $\bR\ge\lambda$ and
$\bR\ge\mu$.  The surviving indices are exactly $\ES(\lambda,\mu,\nu;m)$, which
is the display.

Nonvanishing of the summands, under Hypothesis~\ref{hyp:L}: $\gamma^{\nu}_{\bR}\ne0$
by Lemma~\ref{lem:gnz}; $\overline{\Psi_\chi(\bR)}\ne0$ by admissibility; and
$\overline{\phi_\lambda(\bR)}=\At_{\bR,\lambda}|_{S}\ne0$ by
Hypothesis~\ref{hyp:L} applied with $\bA=\bR$ (which omits $m$) and
$\bB=\lambda$, since $\bR\ge\lambda$ --- and likewise for $\mu$.

(1) Either factor kills the expression; by Lemma~\ref{lem:resfun} a vanishing
residue means $\ell_\nu(f)\in\mathcal{O}$.  (2) is the contrapositive of (1)
together with $\ord_S\ell_\nu(f)\ge-1$ from Lemma~\ref{lem:resfun}.  (3) A sum of
one nonzero term is nonzero.
\end{proof}

\begin{rem}[the single gap, named]\label{rem:gap}
The only thing missing from a clean ``$\iff$'' at fixed $\chi$ is whether
$|\ES|\ge2$ nonzero terms can cancel.  This is the exact analogue of the step
Rick needed Theorem~W and $D^\infty$ for in the pair-divisor case, and I have not
reproduced that machinery here.  Two things keep the statement honest.  First,
\S\ref{sec:quant} removes the gap entirely at the cost of a quantifier on $\chi$,
and the resulting statement is unconditional modulo Hypothesis~\ref{hyp:L}.
Second, I measured how often the question is even asked: $|\ES|\ge2$ occurs in
$2/16$ firing triples at $N=2$, $71/185$ at $N=3$, $1156/1886$ at $N=4$,
$13890/18002$ at $N=5$ and $143692/165134$ at $N=6$, with $|\ES|$ as large as
$32$.  So the census below is \emph{not} a plant drawn from inside its own
hypothesis class: it tests cancellation in a majority of its instances, and no
cancellation occurred.
\end{rem}

%=====================================================================
\section{Theorem S$'$: the converse, with the quantifier where it belongs}
\label{sec:quant}

The pole locus of $\hat c^{\nu}_{\lambda\mu}$ depends on $\chi$ --- that is the
content of \texttt{thm:C} of \texttt{2026-10-17-c2} --- so a converse must either
fix $\chi$ or quantify over it.  Rick's pair-divisor theorem is about the
untwisted constants, where no quantifier is needed; on the boundary family the
untwisted constants have \emph{no} poles ($\kappa=\chi(\varepsilon)=1$ for
$\chi=1$), so the quantified form is the right statement, and it is provable.

\begin{thm}[the boundary converse]\label{thm:Sprime}
Assume Hypothesis~\ref{hyp:L} and \textup{(BD)}.  Fix $N$, $X$, $m$,
$\varepsilon^{h}=1$ and $\lambda,\mu,\nu\in\mathcal{P}$.  Then
\[
  \ES(\lambda,\mu,\nu;m)\ne\varnothing
  \iff
  \begin{array}{l}
  \text{there is an $S$-admissible }\chi\text{ with }\chi(\varepsilon)\ne1\text{ for which}\\
  \hat c^{\nu}_{\lambda\mu}\text{ has a (simple) pole at }S_m^{\varepsilon}.
  \end{array}
\]
Moreover, for each $d\ge N-1$ the set of $\chi$ in the family
$\{\sum_{i=0}^{d}a_iv^{i}\}$ that \emph{fail} to produce the pole is contained in
a proper Zariski-closed subset of the coefficient space.
\end{thm}

\begin{proof}
($\Leftarrow$) is Theorem~\ref{thm:S}(2).

($\Rightarrow$)  Work in $k=\mathbb{C}(q)(w_j:j\ne m)$ and let
$\chi(v)=\sum_{i=0}^{d}a_iv^{i}$ with $a_0,\dots,a_d$ indeterminates, $d\ge N-1$.
Such a $\chi$ is $S$-admissible for $a$ outside a proper closed subset (each
$\chi(w_j)$ and $\chi(\varepsilon)$ is a nonzero linear form in $a$).  By
Theorem~\ref{thm:S},
\[
  \operatorname{Res}_\nu(f)
  =(1-\kappa)\sum_{\bR\in\ES}c_{\bR}\prod_{r\in\bR}L_r,
  \qquad
  L_r:=\chi(w_r)=\sum_{i=0}^{d}a_iw_r^{i},
  \quad
  c_{\bR}:=\gamma^{\nu}_{\bR}\,\overline{\phi_\lambda(\bR)\phi_\mu(\bR)} ,
\]
and each $c_{\bR}\ne0$ by Theorem~\ref{thm:S}.  The $N-1$ linear forms
$\{L_r\}_{r\ne m}$ are $k$-linearly independent in $a$: their coefficient matrix
is the $(N-1)\times(d+1)$ Vandermonde block $(w_r^{i})$ with the $w_r$ pairwise
distinct and $d+1\ge N-1$ columns, hence of full row rank.  Extend
$\{L_r\}_{r\ne m}$ to a basis of the space of linear forms in $a$; in the
resulting polynomial coordinates the squarefree monomials
$\prod_{r\in\bR}L_r$, for distinct subsets $\bR\subseteq[N]\sm\{m\}$, are
distinct monomials, hence linearly independent over $k$.  So
$\sum_{\bR\in\ES}c_{\bR}\prod_{r\in\bR}L_r$ is a \emph{nonzero} polynomial in
$a$ whenever $\ES\ne\varnothing$.  Its zero locus is a proper closed subset, as
is $\{\kappa=1\}$ (a hyperplane) and the non-admissibility locus; a point $a$
outside their union gives an $S$-admissible $\chi$ with $\chi(\varepsilon)\ne1$
and $\operatorname{Res}_\nu(f)\ne0$, i.e.\ a pole, simple by
Theorem~\ref{thm:S}(2).
\end{proof}

\begin{cor}[the pole-free renormalisations, per $\nu$]\label{cor:free}
Let $\chi$ be $S$-admissible for every $m$ and every $\varepsilon^{h}=1$.  If
$\chi(\varepsilon)=1$ for all $\varepsilon^{h}=1$ --- condition \textup{(SR)} of
\texttt{2026-10-17-c2} --- then no $\hat c^{\nu}_{\lambda\mu}$ has a boundary pole,
for any $N$, any $(\lambda,\mu,\nu)$.  If $\chi(\varepsilon)\ne1$ for some
$\varepsilon$, then the set of $\nu$ for which a pole at $S_m^{\varepsilon}$ can
occur is the \textbf{up-set} generated by
$\{\bR\cup\{m\}:\bR\ge\lambda,\bR\ge\mu,\ m\notin\bR\}$, and is empty exactly
when no $\bR\subseteq[N]\sm\{m\}$ dominates both $\lambda$ and $\mu$.
\end{cor}

\begin{proof}
The first half is Theorem~\ref{thm:S}(1).  For the second,
$\ES(\lambda,\mu,\nu;m)\ne\varnothing$ is by definition the condition that some
$\bR\cup\{m\}$ with $\bR\ge\lambda,\mu$ satisfies $\bR\cup\{m\}\le\nu$, which is
monotone in $\nu$ by Corollary~\ref{cor:order}(1).
\end{proof}

%=====================================================================
\section{The $A$ half: there is no converse, and why}
\label{sec:Ahalf}

\begin{prop}\label{prop:A}
Fix $i<j$ and let $H=\{w_i=w_j\}$ with local ring $\mathcal{O}_{ij}$.  For
\emph{every} unit $\Psi$ of the $A$-type GKM ring $\mathcal{G}_{ij}$ --- in
particular for $\Psi=\Psi_\chi$ with $\chi$ any $ij$-admissible renormalisation,
and for $\Psi=1$ --- every structure constant of $\{\Psi\phi_{\bI}\}$ lies in
$\mathcal{O}_{ij}$.  Hence there is no renormalisation in this class for which an
$A$-type pole occurs, and the analogue of Definition~\ref{def:ES} for the
$A$-type family is the empty set for all $(\lambda,\mu,\nu)$.
\end{prop}

\begin{proof}
This is Theorem~B of
\texttt{2026-10-18-atype-divisors-and-the-gkm-lattice.tex}, with
\texttt{lem:unit} supplying $\Psi_\chi\in\mathcal{G}_{ij}^{\times}$; the mechanism
is that $\operatorname{span}_{\mathcal{O}}\{\phi_{\bI}\}=\mathcal{G}_{ij}$ is a
\emph{ring} (there the $A$-type collapse multiplier is $\rho^{A}=1$, by the
$u$-symmetry of Lemma~A1 there, equivalently because $D(e_i-e_j)=\varnothing$ in
the orbit classification of \texttt{2026-10-19-typeB} Lemma~\texttt{lem:orbits}),
and $\Psi\mathcal{G}=\mathcal{G}$ for a unit $\Psi$.  A ring has regular
structure constants.
\end{proof}

\begin{rem}[this is not the same shape of statement as Theorem~\ref{thm:Sprime}]
The asymmetry is structural and worth stating plainly.  On $H_{e_i-e_j}$ the two
colliding fixed points $\bM$ and $\bM\triangle\{i,j\}$ carry the \emph{same
multiset of $w$-values}, so $\Psi_\chi$ --- indeed any element of
$\mathcal{G}_{ij}$ --- has collision multiplier $1$ \emph{for every} $\chi$,
and the lattice is a ring on the nose.  On $S_m^{\varepsilon}$ the value
multisets differ by the single element $w_m$, and the multiplier of $\Psi_\chi$ is
$\chi(\varepsilon)$ --- which is $1$ for the natural $\chi$ but not for all of
them.  The $A$ family therefore has no free parameter left to break, while the
$S$ family has exactly one, and it is $\chi(\varepsilon)$.  In the language of
\texttt{lem:orbits}: $D(e_i-e_j)=\varnothing$ against $D(e_m)=\{m\}$.  The
grouping ``S/A'' in the work assignment collects a family with one degree of
freedom together with a family that has none.
\end{rem}

%=====================================================================
\section{Verification}
\label{sec:verif}

All computation is pure SymPy plus an exact truncated-Laurent-series linear
algebra module over $\mathbb{Q}((t))$ written for this session
(\texttt{proofs/code-1019c2-sa-converse/}).  Orders of vanishing are computed by
Gaussian elimination with minimal-valuation pivoting on the substitution
$w_m=\varepsilon+t$, with $q$ and the remaining $w_j$ fixed rationals; the
arithmetic is over $\mathbb{Q}$ with \texttt{Fraction}, so the only
approximation is the truncation order, and the routines assert they never needed
a coefficient beyond it.  Every run \textbf{states its domain} and \textbf{gates
on genericity} in the sense of Remark~\ref{rem:generic}.

\subsection*{V1 --- the structure-constant census (end-to-end, direct linear solve)}
Domain: $X=C$ at $N=2$ and $N=3$ and $X=B$ at $N=3$; every $m\in[N]$; every
admissible $\varepsilon$ (both signs in type $C$, $\varepsilon=1$ in type $B$);
twists $\chi\in\{1,\ \psi,\ v\psi(v),\ v,\ v+2\}$; and \textbf{all} $2^{3N}$
triples $(\lambda,\mu,\nu)$, no sampling.  $\hat c^{\nu}_{\lambda\mu}$ is obtained
by solving $\Phi x=\Psi_\chi\phi_\lambda\phi_\mu$ --- i.e.\ \emph{without} using
Lemma~\ref{lem:inverse} or any part of \S\ref{sec:res}.

\begin{center}
\begin{tabular}{llrrrr}
\toprule
run & arm & poles & TP & FP & FN\\
\midrule
$N=2$, $C$, $m=1,2$, $\varepsilon=\pm1$ & $\kappa=1$ (3 of 5 twists) & $0$ & --- & $0$ & $0$\\
$N=2$, $C$, $m=1$, $\varepsilon=-1$ & $\kappa=-1$ ($v\psi$, $v$) & $11$ & $11$ & $0$ & $0$\\
$N=2$, $C$, $m=2$, $\varepsilon=-1$ & $\kappa=-1$ & $5$ & $5$ & $0$ & $0$\\
$N=2$, $C$, $m=1$, $\varepsilon=+1$ & $\kappa=3$ ($v+2$) & $11$ & $11$ & $0$ & $0$\\
$N=2$, $C$, $m=2$, $\varepsilon=+1$ & $\kappa=3$ & $5$ & $5$ & $0$ & $0$\\
$N=3$, $C$, $m=1$ & $\kappa\ne1$ & $93$ & $93$ & $0$ & $0$\\
$N=3$, $C$, $m=2$ & $\kappa\ne1$ & $62$ & $62$ & $0$ & $0$\\
$N=3$, $C$, $m=3$ & $\kappa\ne1$ & $30$ & $30$ & $0$ & $0$\\
$N=3$, $C$ and $B$, all $m$ & $\kappa=1$ & $0$ & --- & $0$ & $0$\\
\bottomrule
\end{tabular}
\end{center}

\noindent
Every measured pole order was $-1$ (simple), as Theorem~\ref{thm:S}(2) requires;
no order $\le-2$ occurred anywhere.  The $\kappa=1$ rows are the negative
control: the same twist $v\psi(v)$ that produces $93$ poles at $\varepsilon=-1$
produces $0$ at $\varepsilon=+1$, and $v+2$ does the reverse --- so the
instrument is not simply reporting ``pole'' whenever $\ES\ne\varnothing$.

\subsection*{V2 --- the structural lemmas}
\begin{itemize}
\item \textbf{Corollary~\ref{cor:order}(3)}, exhaustive: $21844/21844$ pairs
  $(\bM,\nu)$, all $N\le7$, $0$ failures.
\item \textbf{Hypothesis~\ref{hyp:L}}: $699/699$ instances with $\bA\ge\bB$,
  $m\notin\bA$; domain $N\le4$, $X\in\{B,C\}$, every $m$, every admissible
  $\varepsilon$, gated-generic point.  \textbf{Plant}: the same quantity for
  $\bA\not\ge\bB$ must vanish by (M3) --- $1179/1179$.  The plant is the positive
  control that the instrument can see a zero at all.
\item \textbf{Lemma~\ref{lem:inverse}} against a numerically inverted $\Phi$:
  $448/448$ entries, $N=2,3$, $X=C$, every $m$, both $\varepsilon$, every $\nu$.
\item \textbf{Lemma~\ref{lem:gamma}} ($\gamma_{\bR\cup\{m\}}+\gamma_{\bR}=0$):
  $224/224$, same domain.
\item \textbf{(BD)} in type $C$: $v_S(\At_{\bI,\bI})=[m\in\bI]$ for every $\bI$,
  by \emph{two} mechanisms (the symmetrisation quotient and the independent
  positive-root product \texttt{diag\_rootproduct}), $N=2,3$, all $m$, both
  $\varepsilon$: agreement on every $\bI$.  In type $B$: $34/34$, $N\le3$.
\end{itemize}

\subsection*{V3 --- the gate: Motegi's all-cardinality biorthogonality}
Domain: $N=1,2,3$; $X\in\{B,C\}$; two random rational $(q,w)$ points each; all
$4^{N}$ pairs $(\bM,\mathbf L)$.  Result: $\mathbf{336/336}$, $0$ mismatches.
The check evaluates $\At$ from the symmetrisation formula, so it is independent
of Motegi's branching-plus-RTV proof.

\subsection*{V4 --- the non-cancellation step, isolated}
Domain: $N=3$ and $N=4$; $X=C$; every $m$; both
$\varepsilon$; all triples; four twists, including one in which $\chi(w_j)$ are
\emph{independently prescribed} rationals --- the computational counterpart of
the interpolation in Theorem~\ref{thm:Sprime}.  The residue sum of
Theorem~\ref{thm:S} is evaluated term by term in $k$ and tested against
$\ES\ne\varnothing$.  $N=3$: $185$ firing triples per twist, $0$ FP, $0$ FN, for
each of four twists --- and $71$ of the $185$ have $|\ES|\ge2$ (up to $|\ES|=4$),
so cancellation was genuinely tested and did not occur.  $N=4$, \textbf{complete}:
$1886$ firing triples per twist summed over $m$ ($736,562,393,195$ for
$m=1,2,3,4$), five twist/$\varepsilon$ arms with $\kappa\ne1$ per $m$, hence
$\mathbf{9430/9430}$ true positives, $0$ FP, $0$ FN --- of which $1156$ of the
$1886$ triples have $|\ES|\ge2$, up to $|\ES|=8$.
\textbf{Provenance, stated:} V4 \emph{uses} Lemma~\ref{lem:inverse}, so it is not
an independent check of that lemma --- it is an independent check of the
\emph{non-cancellation step}, conditional on Lemma~\ref{lem:inverse}, which V2
and V3 verify separately.

\subsection*{Two instrument faults, reported}
\begin{enumerate}
\item \textbf{A non-generic sample read as a phenomenon.}  The first $N=3$ census
  reported $\ord_S\det\Phi=8$ against the predicted $4$.  The random point had
  $w_3=7/3$ with $q=3/7$, i.e.\ $w_3=q^{-1}$, so the root factors
  $(1-qw_mw_3)$ and $(1-qw_3/w_m)$ vanish at $w_m=+1$; the extra order appeared
  for exactly the four $\bI$ predicted by that diagnosis.  The control caught it
  because the predicted number $\#\{\bI\ni m\}$ had been written down
  \emph{before} the run.  Fixed by the genericity gate of
  Remark~\ref{rem:generic}.
\item \textbf{A simplifier that guesses.}  The first biorthogonality run reported
  $20$ mismatches at $N=3$, all containing factors like $2^{41/58}$.  Cause:
  \texttt{sympy.nsimplify} applied to a rational number, which searches for a
  ``nice'' closed form and returns algebraic nonsense.  \texttt{nsimplify} is not
  a simplifier.  Replaced by \texttt{Rational(cancel($\cdot$))}: $336/336$.
  Both faults produced confident, legible, wrong numbers; one read as a
  catastrophe and one as a phenomenon, and the sign of the error said nothing
  about its cause.
\end{enumerate}

%=====================================================================
\section{Gaps, graded}
\label{sec:gaps}

\begin{description}
\item[G1 --- Hypothesis~\ref{hyp:L} ($L^{S}$): \textbf{OPEN}, measured.]
  $\bA\ge\bB$, $m\notin\bA$ $\Rightarrow$ $\At_{\bA,\bB}|_{S}\ne0$.  Needed for
  Lemma~\ref{lem:gnz}, for the ``every summand is nonzero'' clause of
  Theorem~\ref{thm:S}, and hence for Theorem~\ref{thm:Sprime}; \emph{not} needed
  for Theorem~\ref{thm:S}(1) or (2).  $699/699$ with a $1179/1179$ plant.  The
  natural proof is Rick's Lemma~L degeneration $w_c=T^{-e_c}$ at $q=0$; it does
  not obviously survive pinning $w_m=\varepsilon$, and I did not find a
  replacement.  \textbf{This is the first thing to attempt next.}
\item[G2 --- non-cancellation at fixed $\chi$: \textbf{OPEN}, measured, and
  bypassed by a quantifier.]  Whether $|\ES|\ge2$ nonzero summands can cancel for
  a \emph{given} $\chi$.  Theorem~\ref{thm:Sprime} removes the gap by quantifying
  over $\chi$; at fixed $\chi$ the statement remains a one-way implication plus
  the $|\ES|=1$ case.  Measured: $0$ cancellations in $185$ opportunities at $N=3$ (four twists) and
  $0$ in $1886$ at $N=4$ (five twist/$\varepsilon$ arms, $9430$ instances), with
  $|\ES|\ge2$ in $38\%$ and $61\%$ of them respectively.
\item[G3 --- (BD) in type $B$: \textbf{measured, not proved}.]
  $\ord_S\At^{B}_{\bI,\bI}=[m\in\bI]$ on $\{w_m=1\}$.  Type $C$ is proved inside
  \texttt{prop:short} of \texttt{2026-10-17-c2}; the type-$B$ diagonal is
  \texttt{normalizationWB} (L1386) and the needed count is $34/34$ at $N\le3$.
  Everything in \S\S\ref{sec:res}--\ref{sec:quant} is therefore
  \textbf{proved in type $C$ and conditional on G3 in type $B$}.
\item[G4 --- the $\varepsilon=-1$ component in type $B$: \textbf{does not exist},
  not unproved.]  $\{w_m=-1\}$ is not a root hyperplane of $B_N$ and (BC) fails
  there (\texttt{2026-10-19-typeB} \S``beyond'', $0/9$ at $r\ge1$ against $6/6$
  at $r=0$).  Recorded so that the asymmetry between the two types is not later
  read as a missing case.
\item[G5 --- Motegi's biorthogonality is cited, and now corroborated, not
  re-derived.]  Lemma~\ref{lem:inverse} rests on
  \texttt{thm:BC-all-cardinality-biorthogonality}.  I have read the statement and
  its setup first-hand (L4146--4411) and the start of its proof
  (\texttt{lem:allcard-Phihat-inverse}, L4424--4440), and verified the statement
  numerically against my own engine at $N\le3$ ($336/336$).  I have \emph{not}
  read the RTV stable-envelope uniqueness argument it ultimately rests on.  This
  is the same dependency that sits under Rick's $D^\infty$ and his Theorem~W det
  lemma; it is now the one shared, unverified root of three results and should be
  graded as such in both registries.
\end{description}

%=====================================================================
\section{Pre-registration: which arm fired}
\label{sec:prereg}

The brief pre-registered three questions with an explicit third,
\emph{not-well-posed} arm, because on 2026-10-19 cycle~1 all three binary
registers returned an answer that was not on the ballot.  Reported in full,
including arm 3.

\begin{description}
\item[P1 --- ``the S/A converse holds at full $q$ for all $N$''.]
  \textbf{SPLIT: arm 3 for the $A$ half, arm 1 for the $S$ half.}  The $A$ family
  carries no converse because it carries no poles under any renormalisation
  (Proposition~\ref{prop:A}); the $S$ family does, and it is
  Theorem~\ref{thm:Sprime}.  \textbf{The premise of the question --- that ``S/A''
  is one family with one answer --- is what failed.}  The arm-3 detector the
  brief specified (check each candidate divisor is a root hyperplane) fired as
  predicted, but on the $S$ half and in type $B$ only: $\{w_m=-1\}$ is not a root
  hyperplane of $B_N$ (G4).
\item[P2 --- ``$\widehat\Phi$ specialises to my $\At$-restriction at type $A$''.]
  \textbf{Arm 3: $\widehat\Phi$ is a different object} --- the type-$A$,
  fixed-cardinality, unit-diagonal factor of $\mathsf A^{X}$
  (Remark~\ref{rem:gate}).  The gate's purpose was served instead by the parent
  theorem, whose matrix \emph{is} mine up to a diagonal row scaling, because
  $\Delta^{X}_{\bA}=\At_{\bA,\bA}$ by definition.  Verified $336/336$.
\item[P3 --- ``the colength $2^{N-2}$ falls out of biorthogonality in one step,
  canonically''.]  \textbf{Arm 3: ``canonically'' has no content here}, exactly as
  the brief anticipated.  The boundary colength does fall out in one step ---
  $\ord_S\det\Phi=\sum_{\bI}\ord_S\At_{\bI,\bI}=\#\{\bI\ni m\}$, inside the proof
  of Proposition~\ref{prop:lat} --- but that is the same count by the same index
  set, so there is no second index set for the first to be canonically isomorphic
  to.  What biorthogonality actually bought was not the colength but
  Lemma~\ref{lem:gnz}: the \emph{non}-vanishing of $\gamma$ on its support, which
  no colength argument gives.
\end{description}

\begin{rem}[a prediction recorded for next session]
Theorem~\ref{thm:S} makes $\ES$ the \emph{only} combinatorial input to the
boundary residue.  It is therefore falsifiable in a sharp way: if
Hypothesis~\ref{hyp:L} fails for some $(\bA,\bB,m,\varepsilon)$ then the
corresponding $\gamma^{\nu}_{\bR}$ or $\overline{\phi_\lambda(\bR)}$ vanishes and
a triple with $\ES\ne\varnothing$ will show \emph{no} pole for \emph{every}
$\chi$.  So a single FN in V1 or V4 at fixed $\chi$ that survives a change of
$\chi$ refutes $L^{S}$, while an FN that moves with $\chi$ is G2 (cancellation).
\textbf{The two open gaps are distinguishable by one experiment}, and that
experiment is the cheapest next step after attempting $L^{S}$ directly.
\end{rem}

\section*{Questions opened}

None.  The allocator \texttt{memory/questions/NEXT-FREE.md} was not advanced by
this session: G1 and G2 are gaps in a result, not new questions, and
Q530 was answered as part of gate 0b (Remark~\ref{rem:gate}, P3) rather than
re-opened.

\end{document}
